\begin{proposition}[DP37: the true dual torus and a primitive normalized kernel]\label{DP37}
In the actual model of DP32 over $\mathcal O=W(\F_q)$, the automorphism
\[
 A_\zeta(\mathbf X)=(\zeta^pX_1,\ldots,\zeta^{p^{h-1}}X_{h-1}),\qquad\zeta^{q-1}=1,
\]
satisfies $A_\zeta^D=[\zeta]_{G_h}$. Consequently the primitive formal Poincar\'e biextension has an integral normalized frame invariant under
\[
 ([\zeta],A_\zeta^{-1}).
\]
In that frame there is a unique invariant rational logarithmic splitting $\beta$. It vanishes on both axes, has no bidegree-$(1,1)$ term, and has the integral exponential coboundaries required in DP36. Thus for odd $p$ DP36 applies to an actual primitive formal Poincar\'e representative. Its compatible splitting still has the denominators proved in DP35. This does not identify the particular section, gauge, path or harmonic operators of a global theta correspondence.
\end{proposition}
\begin{proof}
The explicit monomial map $M$ of DP32 commutes with $A_\zeta$: the ordinary arrows use $p\cdot p^{i-1}=p^i$, and the exceptional arrow uses $p^{h+1}\equiv p\pmod{q-1}$. Hence its logarithm and integral group law commute as well. This proves that $A_\zeta$ is an integral automorphism. It suffices first to take $\zeta$ primitive. Put $a=A_\zeta^D$ and write $\lambda=da$ on the one-dimensional tangent space of $G_h$.

Tate's functorial Hodge--Tate sequence \cite[\S4, Corollary 2 following Theorem 3, p.~180]{Tate67}, over the completed algebraic closure $C$ of $\operatorname{Frac}(\mathcal O)$, is
\[
 0\longrightarrow \Lie(H)_C
 \longrightarrow \Hom_{\Z_p}(T_pG_h,C)
 \longrightarrow \Lie(G_h)_C^*\otimes_C\Hom_{\Z_p}(\Z_p(1),C)
 \longrightarrow0.
\]
On its middle term $a$ acts by precomposition, on the first term by $d(a^D)=dA_\zeta$, and on the cotangent factor by pullback, with scalar $\lambda$. These are adjoint actions, so neither Hom nor the cotangent introduces an inverse eigenvalue. This also follows directly by differentiating the natural Cartier identity $\langle aP,Q\rangle=\langle P,a^DQ\rangle$. The cyclotomic factor is unchanged.

The characteristic polynomial $P_a(Z)$ of $a$ on the rank-$h$ $\Z_p$ Tate module is consequently
\[
 P_a(Z)=(Z-\lambda)\prod_{i=1}^{h-1}(Z-\zeta^{p^i}).
\]
It is monic of degree $h$ over $\Q_p$. The primitive root $\zeta^p$ has degree exactly $h$: its Frobenius conjugates are $\zeta^{p^j}$, and $p^m\not\equiv1\pmod{p^h-1}$ for $0<m<h$. Therefore its degree-$h$ minimal polynomial divides and equals $P_a$. Comparing the one missing root proves $\lambda=\zeta$. In characteristic zero a one-dimensional formal endomorphism is determined by its differential: conjugate by the normalized logarithm and compare the coefficient of $X^{n-1}Y$ in additivity to kill all degrees $n\ge2$. Thus $a=[\zeta]$. Taking powers proves the equality for every root, and its integral identity persists under base change.

The evaluation pairing giving the primitive Poincar\'e biextension is preserved by $([\zeta],A_\zeta^{-1})$. Functoriality of the pairing construction, not just equality of isomorphism classes, gives a compatible action of the finite group $\Gamma=\mu_{q-1}$ on the biextension and its axial rigidifications. Choose an integral normalized frame $e$ as in DP35 and write $g(e)=c_g e$. The units $c_g$ are one on both axes and satisfy $c_{\sigma g}=\sigma(c_g)c_\sigma$. Set $N=q-1$ and
\[
 u=\left(\prod_{g\in\Gamma}c_g\right)^{1/N}.
\]
The unique root with constant term one is integral: coefficientwise recursion divides only by the unit $N$. It is one on both axes. The cocycle identity gives
\[
 \sigma(u)=u/c_\sigma,
\]
so $e'=ue$ is an invariant integral normalized frame. Its two addition cocycles are invariant integral units.

Use DP35 to obtain a rational logarithmic splitting $\beta_0$ of these cocycles in the additive logarithm coordinates, and average
\[
 \beta=\frac1N\sum_{g\in\Gamma}g\beta_0.
\]
Both additive coboundaries remain the logarithms of the same invariant integral cocycles. The axes remain zero. The characters are exactly $(\zeta,\zeta^{-p},\ldots,\zeta^{-p^{h-1}})$ by the first part of the proof. No bilinear monomial $xy_i$ is invariant, because $1-p^i$ is not divisible by $q-1$ for $1\le i<h$. Thus $\beta$ has no bidegree-$(1,1)$ term. Any two rational splittings differ by a bilinear form by DP35, so the same character observation proves uniqueness among invariant splittings. All hypotheses of DP36 have now been obtained for the primitive class. Its nontriviality and denominator conclusion are unchanged.
\end{proof}

\begin{theorem}[DP38: the first primitive Poincar\'e coefficients]\label{DP38}
For every $p$ and $h\ge2$, over $W(\F_q)$ there are integral Teichm\"uller coordinates $S,Y_1,\ldots,Y_d$ on the Honda line and its actual dual, $d=h-1$, and a normalized integral Poincar\'e frame in which every rational compatible splitting satisfies
\[
 [S^pY_1]B=\frac1p,\qquad [SY_d^p]B=-\frac1p.
\]
These coefficients are unchanged by a rational bilinear gauge and change only by integral amounts under an axes-normalized integral unit frame. The frame is invariant for the true action $(S,Y_i)\mapsto(\zeta S,\zeta^{-p^i}Y_i)$. Its unique invariant logarithmic splitting has no positive term below degree $p+1$, and
\begin{equation}\label{eq:primitive-first}
 (\log B)_{p+1}=B_{p+1}=\frac{S^pY_1-SY_d^p}{p}.
\end{equation}
The same first homogeneous term holds in the normalized logarithm coordinates. Consequently
\[
 D_G^pD_1 B(0)=(p-1)!,\qquad D_GD_d^p B(0)=-(p-1)!.
\]
Both values are units, although both operators are divisible by $p$ on integral formal functions. In the group coordinates,
\[
 w_t(B)=-\infty\quad\left(0<t<\frac1{p+1},\ t\in\Q\right),\qquad
 \liminf_{|\alpha|\to\infty}\frac{v_p(B_\alpha)}{|\alpha|}\le-\frac1{p+1}.
\]
For odd $p$, $h>2$, and any integral $U$ with unit constant term, the last-dual logarithmic path $x=t$, $y_d=\epsilon t$, $y_i=0$ for $i<d$, gives
\begin{equation}\label{eq:primitive-path-unit}
 v_p\left((p+1)![t^{p+1}]\,
 U(E_G(t),E_H(\epsilon te_d))
 \exp\left(\tfrac12\log B(E_G(t),E_H(\epsilon te_d))\right)\right)=0
\end{equation}
for every unit $\epsilon$. This is an exact finite-order assertion, not an asymptotic upper bound on every selected theta jet or a shorter-chain impossibility theorem.
\end{theorem}
\begin{proof}
We give an explicit display calculation, working compatibly over $\mathcal O/p^n$, where $\mathcal O=W(\F_q)$, and then passing to the inverse limit. Let $P$ be free on $e_0,\ldots,e_d$ over the Witt ring, with normal decomposition
\[
 T=We_0,\qquad L=\bigoplus_{j=1}^dWe_j,\qquad Q=L\oplus IT.
\]
With $e_h=e_0$, specify the Frobenius-semilinear operators by
\[
 Fe_0=e_1,\qquad V^{-1}e_j=e_{j+1},\qquad Fe_j=pe_{j+1}\quad(1\le j\le d).
\]
On $IT$ the relation is $V^{-1}(V\xi\,e_0)=\xi Fe_0$. The linearization of $V^{-1}\oplus F$ on $L\oplus T$ is a cyclic permutation and hence an isomorphism. The actual dual display has
\[
 L'=We_0^*,\quad T'=\bigoplus_{j=1}^dWe_j^*,\quad
 V'^{-1}e_0^*=e_1^*,\quad F'e_0^*=pe_1^*,\quad
 F'e_j^*=e_{j+1}^*\ (j\ge1).
\]
These formulas satisfy the pairing identities of Zink, Definition~19 and (22), p.~143. The $h$-fold Frobenius cycles send a basis vector to $p^{h-1}$ times itself in the first display and to $p$ times itself in the dual, with the corresponding Frobenius on coefficients. Both nilpotence conditions hold for $h\ge2$.

By Zink's Theorem~81 and (145)--(146), p.~204, the formal groups are the cokernels of $V^{-1}-\id$. The maps represented by
\[
 \sigma(S)=[S]e_0,\qquad \tau(\mathbf Y)=\sum_{j=1}^d[Y_j]e_j^*
\]
give integral formal coordinates: their tangent maps to $P/Q$ and $P'/Q'$ are identities, so the formal inverse function theorem, proved coefficientwise by its unit linear term, applies. On the special fibre $p[S]=V[S^p]$, and the display relations carry this around the cycle to $[S^q]e_0$. Thus $[p]_G(S)\bmod p=S^q$. The linear coefficient is $p$, so Hazewinkel's integral Lubin--Tate uniqueness theorem, \S5, p.~77 identifies this line with $G_h$. The second display is its actual Cartier dual. No equality of its entire coordinate logarithm with DP32's chosen formula is assumed in the calculation that follows.

Use the evaluation pairing of the displays. The scalar-one section represented by $(1,\sigma(S),\tau(\mathbf Y))$ in Zink's Proposition~106, p.~226 gives a normalized integral frame. If $\kappa_G,\kappa_H$ are the defects of these lifts, defined by
\[
 (V^{-1}-\id)\kappa_G=\sigma(S)+\sigma(S')-\sigma(F_G(S,S'))
\]
and its dual analogue, its cocycles from (218), p.~239 are
\begin{equation}\label{eq:display-cocycles}
 C_G=\operatorname{hex}\langle V^{-1}\kappa_G,\tau(\mathbf Y)\rangle,
 \qquad C_H=\operatorname{hex}\langle\sigma(S),\kappa_H\rangle^{-1}.
\end{equation}
Here $\operatorname{hex}$ is the Artin--Hasse map (210)--(211), p.~236. The inverse is the minus sign in Zink's additively written multiplicative group. Indeed, the quotient relation in Proposition~106 replaces a lift differing by a defect with its scalar multiplied by the corresponding pairing. Thus these signs agree with $C_G=B(F_G)/(BB)$, and similarly for $C_H$. DP35 supplies rational compatible splittings.

Put
\[
 C_p(S,S')=\frac{(S+S')^p-S^p-S'^p}{p},\qquad
 [S]+[S']-[S+S']=V\xi.
\]
The first two ghost polynomials give $\xi_0=-C_p(S,S')$. The integral Witt addition polynomials are weighted homogeneous, so $\xi_i$ has ordinary degree $p^{i+1}$, and every component of $F^r\xi$ has degree at least $p^{r+1}$. This assertion uses this homogeneous carry vector; it is not an assertion about Frobenius and arbitrary coefficient-degree filtrations. The explicit element
\[
 \kappa_0=-V\xi\,e_0-\sum_{j=1}^dF^{j-1}\xi\,e_j
\]
satisfies the exact telescoping identity
\[
 (V^{-1}-\id)\kappa_0=V\xi\,e_0-F^d\xi\,e_0.
\]
The last term has degree at least $p^{d+1}=q$. Modulo the degree-$q$ ideal the group law is therefore ordinary addition. Moreover the true defect equals $\kappa_0$ to that precision: reduce Zink's exact sequence (146) to this nilpotent test algebra and use its injectivity. No inverse of an uncomputed operator is assumed to preserve degree.

The $e_1$ coefficient of $V^{-1}\kappa_0$ is $-\xi$. Multiplication by $[Y_1]$ sends the $i$th Witt component to $-\xi_iY_1^{p^i}$. In the Artin--Hasse product only the zeroth component contributes linearly in $Y_1$, with coefficient one. Thus
\[
 [Y_1]\log C_G=C_p(S,S')+O((S,S')^q).
\]
Writing $b=\log B$, extraction of $S^{p-1}S'Y_1$ from $\delta_G b$ uses ordinary addition and gives $p[S^pY_1]b=1$. Since $b$ has positive degree on both sides, exponentiation cannot alter its $Y$-linear coefficient. This proves $[S^pY_1]B=1/p$.

For the second coefficient restrict both dual inputs to the last coordinate. Again write $[Y]+[Y']-[Y+Y']=V\xi$. This time
\[
 \kappa'_0=-V\xi\,e_d^*-\xi\,e_0^*,\qquad
 (V'^{-1}-\id)\kappa'_0=V\xi\,e_d^*-F\xi\,e_1^*.
\]
The endpoint has degree at least $p^2$. Addition on this source axis is ordinary to that precision, and the true defect has the same reduction by injectivity. Pair with $[S]e_0$ and use the inverse in \eqref{eq:display-cocycles}. It gives
\[
 [S]\log C_H=-C_p(Y,Y')+O((Y,Y')^{p^2}).
\]
Hence $p[SY_d^p]b=-1$, and exponentiation leaves this $S$-linear coefficient unchanged.

These same addition calculations imply $\log_G(S)=S+O(S^q)$ and
\[
 \log_H(Ye_d)=Ye_d+O(Y^{p^2}):
\]
compare the low-degree terms of the logarithm's additive identity on the indicated inputs. Therefore a rational bilinear change $\log_G(S)\sum_j a_j\log_{H,j}(\mathbf Y)$ contributes to neither of the two coefficients. An axes-normalized unit frame change adds or subtracts $\log I$. Its $Y$-linear and $S$-linear coefficients are integral, since higher powers of $I-1$ cannot contribute there. This proves the claimed invariance of their valuations.

For the full first homogeneous part, the diagonal display map $e_j\mapsto[\zeta^{p^j}]e_j$ commutes with the operators, including Frobenius on its coefficients. It sends $\sigma(S)$ to $\sigma(\zeta S)$, so gives the true $[\zeta]$ on the line. Its inverse dual sends $\tau(\mathbf Y)$ to $\tau(\zeta^{-p^j}Y_j)$. The evaluation pairing and the scalar-one quotient frame are invariant. As in DP37, averaging gives the unique invariant rational $b$.

Suppose an invariant mixed monomial has side degrees $a,b\ge1$ and total degree at most $p+1$. Then
\[
 a-\sum_jp^jb_j\equiv0\pmod{q-1},\qquad
 pb\le\sum_jp^jb_j\le p^db\le q.
\]
If the sum equals $a$, the degree bound forces $b=1$, $a=p$, and $j=1$. If the sum exceeds $a$, their difference can only be $q-1$; the upper bound forces $a=1$, $b=p$, and all dual factors to be $Y_d$. These are precisely the two monomials already computed, both of degree $p+1$. This proves \eqref{eq:primitive-first} and absence of smaller terms. Substitution of the inverse logarithms, whose linear parts are identities, preserves that first homogeneous part. The two stated factorial mixed derivatives follow.

For an integral function, the translation coefficient at source degree $p$ on the line is $D_G^p/p!$, since its logarithm has no nonlinear term of that degree. The corresponding last-dual coefficient is $D_d^p/p!$ by the last-axis calculation above. Both translation operators are integral, so $D_G^p$ and $D_d^p$ are divisible by $p$ on the entire integral function ring. This proves the asserted contrast with the two unit mixed jets. DP35 applied to a degree-$p+1$ coefficient of valuation $-1$ gives the numerical denominator rate.

Finally assume $h>2$. On the stated logarithmic path the first monomial in \eqref{eq:primitive-first} vanishes because $y_1=0$, leaving $-\epsilon^pt^{p+1}/p$. Put $D=D_G+\epsilon D_d$. Its two pure $p$th powers and every interior binomial term are divisible by $p$, so $D^p$ takes integral functions into $p\mathcal O[[S,\mathbf Y]]$. Consequently the integral $U$ factor contributes a multiple of $p$ to its $(p+1)$st derivative. There is no earlier positive term in the exponential gauge; its contribution at that order is
\[
 -U(0)\epsilon^p\frac{(p+1)!}{2p},
\]
a unit for odd $p$. It cannot cancel the first contribution. This proves \eqref{eq:primitive-path-unit}.
\end{proof}

At $h=2$ the path coefficient is $(\epsilon-\epsilon^p)/p$ and can vanish, which explains the restriction in \eqref{eq:primitive-path-unit}. For any prime the same unit-jet proof works with $\exp(k\log B)$ for a unit scalar $k$ in place of the half-power. The asymptotic denominator rate also survives on a formal generic group-coordinate diagonal after an unramified extension with transcendental residue unit $\epsilon$: first restrict to $S,Y_1$, which retains the integral first coboundary and the coefficient $1/p$, then put $S=t,Y_1=\epsilon t$. In each nonzero homogeneous part the least-valuation residue is a nonzero polynomial in $\bar\epsilon$, so no cancellation of that minimum occurs. This proves the same rate there, but makes no claim for a prescribed algebraic CM parameter or for an arbitrary logarithmic diagonal.

\begin{proposition}[DP39: what harmonicity removes, and what it does not]\label{DP39}
The following two assertions distinguish full-series denominators from selected theta-type jets.
\begin{enumerate}[leftmargin=*,label=(\alph*)]
\item There are integral formal-biextension data, two canonical gauges, and a primitive germ $f=1$ for which one normalized series has unbounded denominators, while all the specified harmonic values agree. In fact the agreement holds for every germ $f$ whenever the polynomial $P$ is killed by the specified second-order contraction.
\item For odd $p$ and every $h\ge2$, two copies of the actual primitive pair have an integral polynomial satisfying the summed harmonic, row-invariance and fixed-row-degree conditions defining \path{HaO}, whose normalized kernel value is a unit. The same invariant polynomial takes values in $pR$ on every integral function on that product. Therefore those algebraic harmonic conditions alone do not transfer the bare integral-function divisibility through the primitive exponential.
\end{enumerate}
Neither assertion constructs a new global mixed-type theta correspondence or asserts that the particular germ $1$ is a global theta section.
\end{proposition}
\begin{proof}
For a polynomial $P$ and a formal series $f$ over a characteristic-zero field, write $\langle P,f\rangle=P(\partial)f(0)$. For every polynomial $Q$,
\begin{equation}\label{eq:fischer-ideal}
 \langle P,Qf\rangle=\langle Q(\partial)P,f\rangle.
\end{equation}
For monomials $P=Z^\alpha,Q=Z^\beta,f=Z^\gamma$, both sides vanish unless $\alpha=\beta+\gamma$, when both are $\alpha!$. Linearity and the finite jet read by $P$ prove the identity for formal series. If $Q_j(\partial)P=0$ for every $j$, the pairing kills the formal ideal $(Q_j)$. If $\eta$ has zero constant term and belongs to that ideal, then $\exp\eta-1=\eta\sum_{n\ge1}\eta^{n-1}/n!$ also belongs to it. Hence
\[
 \langle P,f\exp\eta\rangle=\langle P,f\rangle.
\]

For (a), let $F=\widehat\Ga^2$ over $\Z_p$, let $O=\Z\times\Z$ act diagonally, and let its involution exchange the two factors. Take trivial addition, symmetry and adjoint cocycles. Both $h_0=0$ and
\[
 h_1(X,Y)=X_1Y_2+X_2Y_1
\]
are canonical gauges: their axes vanish, their addition coboundaries vanish by biadditivity, and symmetry is immediate. For $(a,b)\in O$ the adjoint identity is explicitly
\[
 h_1((aX_1,bX_2),Y)=aX_1Y_2+bX_2Y_1
 =h_1(X,(bY_1,aY_2)).
\]
Take $w_0=(1,0)$, so $w_0+w_0^*=1$ and $h_{1,L}=X_1X_2$. For $f=1$ the normalized series is $\exp(X_1X_2)$, whose coefficients $1/n!$ have unbounded denominators. If $\partial_U\partial_VP=0$, its selected value agrees with that in gauge zero for every $f$, by \eqref{eq:fischer-ideal} with $Q=X_1X_2$. Equivalently $P$ is a sum of pure polynomials in $U$ and in $V$, and the exponential is one on each axis. The second-order operator in the source's symmetric-matrix convention is $2\partial_U\partial_V$, with the same kernel over $\Q_p$, including when $p=2$. For $f=1$ all positive-degree selected values vanish. This proves (a) and shows why primitive initial terms and full-series nonintegrality do not by themselves decide harmonic integrality.

For (b), use DP38's invariant splitting in logarithm coordinates, whose coefficient is $[x^py_1]B=1/p$. On two copies put
\[
 \Theta=B(x_1,\mathbf y^{(1)})^{1/2}B(x_2,\mathbf y^{(2)})^{1/2},\qquad
 P=X_1^pY_{1,1}-pX_1^{p-1}X_2Y_{1,2}.
\]
Give the positive type one row and the negative type $d$ rows, each with two columns. Give the positive row degree $p$, the first negative row degree one, and all other rows degree zero. Every coefficient is integral and every monomial has those degrees. The positive row has no earlier row. Adding an earlier negative row to a later one leaves $P$ unchanged because it uses only the first negative row. Thus the prescribed row invariances hold. For the identity-form contractions
\[
 \Delta_i=\partial_{X_1}\partial_{Y_{i,1}}+
          \partial_{X_2}\partial_{Y_{i,2}},
\]
one has $\Delta_iP=0$ for $i>1$ and
\[
 \Delta_1P=pX_1^{p-1}-pX_1^{p-1}=0.
\]
These are exactly the algebraic definitions \path{lap}, \path{harmK}, \path{RowInv}, \path{RowDeg}, and \path{HaO} in
\path{Formalism/LocalToGlobal/ProjectiveLattices/RowDeg.lean},
lines 21--30 and 37--48; \path{IsThRule.lean}, \path{harmA}, lines 86--87, specializes this definition in the source. They are sums over columns, not individual contractions.

Every monomial of $B-1$ has positive degree on both sides. Therefore the $y_1$-linear coefficient of $B^{1/2}$ is half that of $B$, and the first monomial of $P$ contributes
\[
 p!\,[x^py_1]B^{1/2}=\frac{(p-1)!}{2},
\]
a unit. The second monomial contributes zero: it differentiates the first copy $p-1>0$ times in $x_1$, with its entire dual variable set zero, where that factor is identically one. Hence $P(\partial)\Theta(0)=(p-1)!/2$.

For an integral function $f$ on $(G_h\times H)^2$, the first term of $P(D)f$ is divisible by $p$ because $D_{x_1}^p$ is coefficientwise divisible by $p$ at height at least two. Its other invariant derivation preserves that ideal. The second term has its explicit factor $p$ and otherwise only integral invariant derivations. Thus $P(D)f(0)\in pR$. This proves (b), with no assumption that the formal row labels already define a new global source datum.
\end{proof}

The associated ideal in \eqref{eq:fischer-ideal} is generated by column sums $\sum_bX_{\tau,a,b}X_{\bar\tau,a',b}$, not every separate product. In DP22's pure CM type $\{1,2,3\}\subset(\Z/7)^\times$, the conjugate type is $\{6,5,4\}$ and is disjoint. If this formula is applied with one row on the type and zero outside it, each conjugate-row contraction has an empty index set. It does not impose $\partial_x\partial_{y_i}$ on the height-three line and its dual directions. This observation is about the stated harmonic formula; DP25 still prevents treating that row list as the old induced-type global host.

